\documentclass[a4paper]{article}
\usepackage[a4paper,margin=0mm]{geometry}
\usepackage{fontspec}
\usepackage{unicode-math}
\usepackage{tikz}
\usetikzlibrary{fit,calc,positioning,arrows.meta}
\usepackage[table]{xcolor}
\usepackage[absolute]{textpos}
\usepackage{array}

\setmainfont{Libertinus Serif}
\setsansfont{Libertinus Sans}
\setmathfont{Libertinus Math}

\definecolor{accent}{HTML}{0F6A52}
\definecolor{badge}{HTML}{A9CFC2}
\definecolor{ink}{HTML}{2A2E30}

\setlength{\TPHorizModule}{1mm}
\setlength{\TPVertModule}{1mm}
\setlength{\parindent}{0pt}
\pagestyle{empty}

\newcommand{\topicheader}[2]{%
\begin{textblock*}{194mm}(16mm,10mm)
\begin{tikzpicture}
\node[fill=accent,text=white,font=\sffamily\bfseries\fontsize{13}{13}\selectfont,
      minimum width=9mm,minimum height=9mm,inner sep=0pt] at (4.5mm,-4.5mm) {#1};
\node[font=\sffamily\bfseries\fontsize{15}{16}\selectfont,text=ink,anchor=west] at (12mm,-4.5mm) {#2};
\end{tikzpicture}
\end{textblock*}
\begin{textblock*}{194mm}(16mm,22.5mm)
\begin{tikzpicture}
\fill[accent] (0,0) rectangle (194mm,0.5mm);
\end{tikzpicture}
\end{textblock*}
}

\newcommand{\pagekicker}[1]{%
\begin{textblock*}{100mm}(16mm,25.5mm)
\begin{tikzpicture}
\node[fill=ink,text=white,font=\sffamily\bfseries\fontsize{8}{8}\selectfont,
      inner sep=2.2pt,rounded corners=1pt] {\MakeUppercase{#1}};
\end{tikzpicture}
\end{textblock*}
}

\newcommand{\videobanner}[1]{%
\begin{textblock*}{178mm}(16mm,25.5mm)
\begin{tikzpicture}
\fill[accent] (0,0) rectangle (178mm,-10mm);
\draw[white,line width=0.6pt] (3mm,-1.7mm) rectangle (11.5mm,-8.3mm);
\node[white,font=\sffamily\fontsize{4}{4.4}\selectfont,align=center] at (7.25mm,-5mm) {QR\\CODE};
\node[white,font=\sffamily\bfseries\fontsize{9.5}{11}\selectfont,anchor=west] at (16mm,-5.6mm) {WATCH THE VIDEO BEFORE YOU SOLVE};
\node[white,font=\sffamily\fontsize{7.5}{8}\selectfont,anchor=east] at (174mm,-5.6mm) {YouTube ~\textbullet~ #1};
\end{tikzpicture}
\end{textblock*}
}

\newcommand{\pagefooter}[1]{%
\begin{textblock*}{194mm}(8mm,286.41mm)
  \begin{tikzpicture}
    \fill[accent] (0,0) rectangle (194mm,0.6mm);
  \end{tikzpicture}
\end{textblock*}
\begin{textblock*}{63.59mm}(8mm,288.27mm)
  {\sffamily\bfseries\color{accent}\fontsize{13}{16}\selectfont SAT 1200+}
\end{textblock*}
\begin{textblock*}{10mm}(188mm,288.27mm)
  \centering{\sffamily\color{accent}\fontsize{14}{16}\selectfont #1}
\end{textblock*}
}

\newcommand{\cardbg}[4]{%
\begin{textblock*}{#3mm}(#1mm,#2mm)
\begin{tikzpicture}
\fill[badge!22,rounded corners=3mm] (0,0) rectangle (#3mm,-#4mm);
\end{tikzpicture}
\end{textblock*}%
}

\newcommand{\pill}[3]{%
\begin{textblock*}{160mm}(#1mm,#2mm)
\begin{tikzpicture}
\node[fill=accent,text=white,anchor=north west,font=\sffamily\bfseries\fontsize{10}{10}\selectfont,
      rounded corners=8pt,inner xsep=5mm,inner ysep=1.8mm] at (0,0) {RULE ~\textperiodcentered~ #3};
\end{tikzpicture}
\end{textblock*}%
}

\newcommand{\pillw}[3]{%
\begin{textblock*}{160mm}(#1mm,#2mm)
\begin{tikzpicture}
\node[fill=accent,text=white,anchor=north west,font=\sffamily\bfseries\fontsize{10}{10}\selectfont,
      rounded corners=8pt,inner xsep=5mm,inner ysep=1.8mm] at (0,0) {#3};
\end{tikzpicture}
\end{textblock*}%
}

\arrayrulecolor{accent}
\setlength{\arrayrulewidth}{0.7pt}
\renewcommand{\arraystretch}{1.5}


\begin{document}
\topicheader{9}{Hard Triangles \& Chained Geometry}
\videobanner{Lesson 9}
\pagefooter{1}
\mbox{}%
% ---- Right-triangle ratios (SOH-CAH-TOA) (y=40.0 h=37.0 end 77.0)
\cardbg{16}{40.0}{178}{37.0}
\pill{26}{43.0}{Right-triangle ratios (SOH-CAH-TOA)}
\begin{textblock*}{158mm}(26mm,51.5mm)
{\fontsize{15}{18}\selectfont $\sin A=\dfrac{\text{opp}}{\text{hyp}}$ \qquad $\cos A=\dfrac{\text{adj}}{\text{hyp}}$ \qquad $\tan A=\dfrac{\text{opp}}{\text{adj}}$}
\end{textblock*}
\begin{textblock*}{158mm}(26mm,64.0mm)
{\fontsize{13}{16}\selectfont legs $5,12$, hyp $13$: \ $\sin A=\dfrac5{13}$, $\cos A=\dfrac{12}{13}$, $\tan A=\dfrac5{12}$}
\end{textblock*}
% ---- Complementary angles (y=83.0 h=34.5 end 117.5)
\cardbg{16}{83.0}{178}{34.5}
\pill{26}{86.0}{Complementary angles}
\begin{textblock*}{158mm}(26mm,94.5mm)
{\fontsize{15}{18}\selectfont $\sin x=\cos(90^\circ-x)$ \qquad $\tan x=\dfrac1{\tan(90^\circ-x)}$}
\end{textblock*}
\begin{textblock*}{158mm}(26mm,107.0mm)
{\fontsize{13}{16}\selectfont $\sin P=\cos Q$ (acute) $\Rightarrow P+Q=90^\circ$ \quad solve for $x$, reject impossible roots}
\end{textblock*}
% ---- From a ratio to the sides: scale with $k$ (y=123.5 h=37.0 end 160.5)
\cardbg{16}{123.5}{178}{37.0}
\pill{26}{126.5}{From a ratio to the sides: scale with $k$}
\begin{textblock*}{158mm}(26mm,135.0mm)
{\fontsize{15}{18}\selectfont $\tan\theta=\dfrac ab\Rightarrow$ legs $ak,\ bk$, hyp $=k\sqrt{a^2+b^2}$}
\end{textblock*}
\begin{textblock*}{158mm}(26mm,147.5mm)
{\fontsize{13}{16}\selectfont $\tan\theta=\dfrac{24}7$, hyp $50$: \ $7k,24k,25k$ with $25k=50$, $k=2$, area $=\dfrac12(14)(48)=336$}
\end{textblock*}
% ---- Special right triangles (y=166.5 h=30.0 end 196.5)
\cardbg{16}{166.5}{178}{30.0}
\pill{26}{169.5}{Special right triangles}
\begin{textblock*}{158mm}(26mm,177.5mm)
{\fontsize{15}{18}\selectfont $45^\circ$-$45^\circ$-$90^\circ$: $\ 1:1:\sqrt2$ \qquad $30^\circ$-$60^\circ$-$90^\circ$: $\ 1:\sqrt3:2$}
\end{textblock*}
\begin{textblock*}{158mm}(26mm,186.0mm)
{\fontsize{13}{16}\selectfont hyp $12$ isosceles right: leg $6\sqrt2$ \qquad $30^\circ$ with hyp $14$: legs $7$ and $7\sqrt3$}
\end{textblock*}
% ---- Pythagorean triples (and their multiples) (y=202.5 h=30.0 end 232.5)
\cardbg{16}{202.5}{178}{30.0}
\pill{26}{205.5}{Pythagorean triples (and their multiples)}
\begin{textblock*}{158mm}(26mm,213.5mm)
{\fontsize{15}{18}\selectfont $3,4,5\quad5,12,13\quad8,15,17\quad7,24,25\quad9,40,41\quad20,21,29$}
\end{textblock*}
\begin{textblock*}{158mm}(26mm,222.0mm)
{\fontsize{13}{16}\selectfont $a^2+b^2=c^2$ \quad $x,\,x+7,\,x+8$ is a right triangle $\Rightarrow x=5$: \ $5,12,13$}
\end{textblock*}
% ---- Altitude to the hypotenuse (y=238.5 h=37.0 end 275.5)
\cardbg{16}{238.5}{178}{37.0}
\pill{26}{241.5}{Altitude to the hypotenuse}
\begin{textblock*}{158mm}(26mm,250.0mm)
{\fontsize{15}{18}\selectfont $h^2=p\,q$ \qquad $a^2=p\,c$ \qquad $b^2=q\,c$ \qquad $h=\dfrac{ab}{c}$ \qquad three similar triangles}
\end{textblock*}
\begin{textblock*}{158mm}(26mm,262.5mm)
{\fontsize{13}{16}\selectfont $p=4$, $q=9$: \ $h=6$, $c=13$, area $39$ \qquad $\angle ABD=\angle C$ (small angles repeat)}
\end{textblock*}
\clearpage
\topicheader{9}{Hard Triangles \& Chained Geometry}
\pagekicker{Rules sheet -- 2/3}
\pagefooter{2}
\mbox{}%
% ---- Similar triangles (y=33.0 h=37.0 end 70.0)
\cardbg{16}{33.0}{178}{37.0}
\pill{26}{36.0}{Similar triangles}
\begin{textblock*}{158mm}(26mm,44.5mm)
{\fontsize{15}{18}\selectfont sides $\times k$ \qquad perimeter $\times k$ \qquad area $\times k^2$ \qquad $DE\parallel BC\Rightarrow\triangle ADE\sim\triangle ABC$}
\end{textblock*}
\begin{textblock*}{158mm}(26mm,57.0mm)
{\fontsize{13}{16}\selectfont $AD=9$, $AB=15$, $DE=12$: \ $BC=12\cdot\dfrac{15}9=20$ \qquad areas $18$ and $98$: $k=\dfrac37$}
\end{textblock*}
% ---- Midsegment and centroid (y=76.0 h=37.0 end 113.0)
\cardbg{16}{76.0}{178}{37.0}
\pill{26}{79.0}{Midsegment and centroid}
\begin{textblock*}{158mm}(26mm,87.5mm)
{\fontsize{15}{18}\selectfont midsegment $=\dfrac12$ third side, parallel to it \qquad centroid splits a median $2:1$}
\end{textblock*}
\begin{textblock*}{158mm}(26mm,100.0mm)
{\fontsize{13}{16}\selectfont $MN=2x+3=\dfrac12(5x-3)\Rightarrow x=9$ \qquad $AG=2\,GD$, so $AD=3\,GD$}
\end{textblock*}
% ---- Triangle angle facts (y=119.0 h=30.0 end 149.0)
\cardbg{16}{119.0}{178}{30.0}
\pill{26}{122.0}{Triangle angle facts}
\begin{textblock*}{158mm}(26mm,130.0mm)
{\fontsize{15}{18}\selectfont sum $180^\circ$ \qquad exterior $=$ far interior angles \qquad isosceles: equal base angles}
\end{textblock*}
\begin{textblock*}{158mm}(26mm,138.5mm)
{\fontsize{13}{16}\selectfont third side: $|a-b|<c<a+b$ \qquad $9,12,14$: $81+144>196$, so the largest angle is \textbf{acute}}
\end{textblock*}
% ---- Parallel lines and polygons (y=155.0 h=32.5 end 187.5)
\cardbg{16}{155.0}{178}{32.5}
\pill{26}{158.0}{Parallel lines and polygons}
\begin{textblock*}{158mm}(26mm,166.0mm)
{\fontsize{15}{18}\selectfont alternate, corresponding: equal \qquad same-side interior: sum $180^\circ$}
\end{textblock*}
\begin{textblock*}{158mm}(26mm,174.5mm)
{\fontsize{13}{16}\selectfont polygon interior sum $(n-2)180^\circ$ \quad regular: exterior $=\dfrac{360^\circ}{n}$, interior $=180^\circ-$exterior}
\end{textblock*}
% ---- Equilateral and isosceles triangles (y=193.5 h=37.0 end 230.5)
\cardbg{16}{193.5}{178}{37.0}
\pill{26}{196.5}{Equilateral and isosceles triangles}
\begin{textblock*}{158mm}(26mm,205.0mm)
{\fontsize{15}{18}\selectfont equilateral: $h=\dfrac{\sqrt3}2\,s$ \qquad area $=\dfrac{\sqrt3}4\,s^2$ \qquad isosceles: $h=\sqrt{\text{leg}^2-\left(\tfrac{\text{base}}2\right)^2}$}
\end{textblock*}
\begin{textblock*}{158mm}(26mm,217.5mm)
{\fontsize{13}{16}\selectfont legs $10$, base $16$: \ $h=6$, area $=48$ \qquad height $15$ equilateral: $s=10\sqrt3$, area $75\sqrt3$}
\end{textblock*}
\clearpage
\topicheader{9}{Hard Triangles \& Chained Geometry}
\pagekicker{Rules sheet -- 3/3}
\pagefooter{3}
\mbox{}%
% ---- Angle of elevation and depression (y=33.0 h=34.5 end 67.5)
\cardbg{16}{33.0}{178}{34.5}
\pill{26}{36.0}{Angle of elevation and depression}
\begin{textblock*}{158mm}(26mm,44.5mm)
{\fontsize{15}{18}\selectfont elevation $=$ depression \qquad two points, one height: $d=h(\cot a-\cot b)$}
\end{textblock*}
\begin{textblock*}{158mm}(26mm,57.0mm)
{\fontsize{13}{16}\selectfont $45^\circ$ and $30^\circ$, $d=40$: \ $h\,(\sqrt3-1)=40$, so $h=20(\sqrt3+1)$}
\end{textblock*}
% ---- Coordinate chains (y=73.5 h=34.5 end 108.0)
\cardbg{16}{73.5}{178}{34.5}
\pill{26}{76.5}{Coordinate chains}
\begin{textblock*}{158mm}(26mm,85.0mm)
{\fontsize{15}{18}\selectfont distance $\sqrt{(\Delta x)^2+(\Delta y)^2}$ \qquad midpoint $\left(\tfrac{x_1+x_2}2,\tfrac{y_1+y_2}2\right)$ \qquad $m_1m_2=-1$}
\end{textblock*}
\begin{textblock*}{158mm}(26mm,97.5mm)
{\fontsize{13}{16}\selectfont right angle at $Q$: $\overrightarrow{QP}\cdot\overrightarrow{QR}=0$ \qquad equidistant from $A,B$: set $PA^2=PB^2$ (linear in $x$)}
\end{textblock*}
% ---- Area toolkit and composites (y=114.0 h=34.5 end 148.5)
\cardbg{16}{114.0}{178}{34.5}
\pill{26}{117.0}{Area toolkit and composites}
\begin{textblock*}{158mm}(26mm,125.5mm)
{\fontsize{15}{18}\selectfont triangle $\dfrac12bh$ \quad trapezoid $\dfrac12(a+b)h$ \quad height $h=s\sin\theta$}
\end{textblock*}
\begin{textblock*}{158mm}(26mm,138.0mm)
{\fontsize{13}{16}\selectfont hidden height first (Pythagoras); \textbf{subtract} corner pieces from a bounding box}
\end{textblock*}
% ---- Chain rule for hard geometry (y=154.5 h=30.0 end 184.5)
\cardbg{16}{154.5}{178}{30.0}
\pill{26}{157.5}{Chain rule for hard geometry}
\begin{textblock*}{158mm}(26mm,165.5mm)
{\fontsize{15}{18}\selectfont name the \textbf{shared} side, angle or $x$ \ $\to$ \ rule 1 \ $\to$ \ rule 2}
\end{textblock*}
\begin{textblock*}{158mm}(26mm,174.0mm)
{\fontsize{13}{16}\selectfont two right triangles, one shared leg: \ $c_1^2-a^2=c_2^2-b^2$ \qquad do not trust a picture}
\end{textblock*}
% ---- New words (y=190.5 h=79.5)
\cardbg{16}{190.5}{178}{79.5}
\pillw{26}{193.5}{NEW WORDS ~\textperiodcentered~ YANGI SO'ZLAR}
\begin{textblock*}{80mm}(26mm,203.0mm)
{\fontsize{10.5}{12}\selectfont
\begin{tabular}{@{}>{\bfseries}p{37mm}p{36mm}@{}}
\cellcolor{accent}\textcolor{white}{\sffamily English} & \cellcolor{accent}\textcolor{white}{\sffamily O'zbekcha}\\
\rowcolor{badge!45} hypotenuse & {\itshape gipotenuza}\\
\rowcolor{white} leg (of a right triangle) & {\itshape katet}\\
\rowcolor{badge!45} altitude (height) & {\itshape balandlik}\\
\rowcolor{white} median & {\itshape mediana}\\
\rowcolor{badge!45} centroid & {\itshape og'irlik markazi}\\
\rowcolor{white} midsegment & {\itshape o'rta chiziq}\\
\rowcolor{badge!45} similar triangles & {\itshape o'xshash uchburchaklar}\\
\rowcolor{white} scale factor & {\itshape o'xshashlik koeffitsiyenti}\\
\end{tabular}}
\end{textblock*}
\begin{textblock*}{80mm}(108mm,203.0mm)
{\fontsize{10.5}{12}\selectfont
\begin{tabular}{@{}>{\bfseries}p{37mm}p{36mm}@{}}
\cellcolor{accent}\textcolor{white}{\sffamily English} & \cellcolor{accent}\textcolor{white}{\sffamily O'zbekcha}\\
\rowcolor{badge!45} exterior angle & {\itshape tashqi burchak}\\
\rowcolor{white} complementary angles & {\itshape to'ldiruvchi burchaklar}\\
\rowcolor{badge!45} transversal & {\itshape kesuvchi to'g'ri chiziq}\\
\rowcolor{white} angle of elevation & {\itshape ko'tarilish burchagi}\\
\rowcolor{badge!45} angle of depression & {\itshape pasayish burchagi}\\
\rowcolor{white} polygon & {\itshape ko'pburchak}\\
\rowcolor{badge!45} perpendicular & {\itshape perpendikulyar}\\
\end{tabular}}
\end{textblock*}
\end{document}
